\documentclass[11pt,a4paper]{article}
\usepackage{fontspec,amsmath,amssymb,amsthm,microtype}
\setmainfont{lmroman10-regular.otf}[BoldFont=lmroman10-bold.otf,ItalicFont=lmroman10-italic.otf,BoldItalicFont=lmroman10-bolditalic.otf]
\usepackage[margin=25mm]{geometry}
\usepackage[unicode,hidelinks]{hyperref}
\setlength{\emergencystretch}{3em}
\setlength{\parskip}{0.4em}
\setlength{\parindent}{0pt}
\allowdisplaybreaks[2]
\raggedbottom
\title{\texorpdfstring{The relative-\allowbreak{}density log-\allowbreak{}Sobolev convention}{The relative-density log-Sobolev convention}}
\author{Omer Deniz Akyildiz\\[0.4em]\normalsize Prepared with AI assistance}
\date{}
\begin{document}\maketitle
Draft mathematical exposition; these notes may contain mistakes.


Write \(\rho=r\nu\), \(H_\nu(\rho)=\int r\log r\,d\nu\), and \(J_\nu(\rho)=\int\|\nabla\log r\|^2d\rho\). This module defines the explicit smooth positive-density LSI convention \(2\alpha H\le J\) and guards real-valued conversions by finite divergence and Fisher information.



A relative density \(r\) is measured against the target \(\nu\), so its induced law is \(\rho_r=r\nu\). For a positive normalized density, \(H_\nu(r)=\int r\log r\,d\nu\) when the integrand is integrable, and \(I_\nu(r)=\int r\|\nabla\log r\|^2d\nu\). The convention \(2\alpha H_\nu(r)\le I_\nu(r)\) fixes the meaning of the LSI constant. These are initially nonnegative extended quantities; finiteness must precede ordinary real integral manipulations.


\section{\texorpdfstring{The law specified by a relative density}{The law specified by a relative density}}


For a real relative density \(r\), define \(\rho_r(dx)=r(x)\nu(dx)\) using a nonnegative scalar \ensuremath{\circ} r). Positivity and normalization are later hypotheses, rather than part of this raw measure definition. A normalized positive \(r\) is the density of the current law relative to the target.



The density specifies a measure before it specifies a probability. To use entropy as a nonnegative divergence, require \(r\ge0\) and \(\int r\,d\nu=1\). Then integration against the induced law is ordinary weighted integration against \(\nu\). Strict positivity is added for logarithmic derivatives, rather than being needed merely to form the weighted measure.


\section{\texorpdfstring{Relative Fisher information}{Relative Fisher information}}


Define \(I_\nu(r)=\int\|\nabla\log r\|^2d\rho_r\) as a nonnegative extended integral. The specified positive density is differentiated pointwise. Subsequent energy estimates establish finiteness before this quantity is used as an ordinary real number.



The relative score differentiates the chosen positive density \(r\) with respect to the reference \(\nu\). For a Gibbs reference with Lebesgue density \(J\), if \(r=p/J\), it is \(\nabla\log p+\nabla U\). Thus Fisher information measures departure from the target score, rather than just roughness of the Lebesgue density.


\section{\texorpdfstring{The log-\allowbreak{}Sobolev inequality}{The log-Sobolev inequality}}


The target is a probability and satisfies \(2\alpha\operatorname{KL}(\rho_r\mid\nu)\le I_\nu(r)\) for every strictly positive \(C^1\) normalized relative density. The factor \(2\alpha\) fixes the convention used throughout the proof. This test class and its normalization are explicit; the definition contains no algorithm or entropy-decay claim.



The hypothesis is a functional inequality for every member of the displayed positive \(C^1\) normalized test class. It is not an assertion that every probability law already has a smooth density. In the ULA application, Gaussian smoothing proves that the actual positive-time relative density belongs to this class before LSI is invoked.


\section{\texorpdfstring{The finite real log-\allowbreak{}Sobolev bound}{The finite real log-Sobolev bound}}


Apply LSI to the positive \(C^1\) normalized density. If KL and Fisher information are finite, this is an inequality between ordinary nonnegative numbers, giving \(2\alpha H_\nu(r)\le I_\nu(r)\) with the same normalization.



The extended inequality does not by itself justify ordinary subtraction or multiplication when infinite values are possible when infinity is possible. Once KL and Fisher are finite, both are ordinary nonnegative real quantities and the coefficient \(2\alpha\) has its usual value. The displayed real LSI then follows without changing its normalization.


\end{document}
